\documentclass[11pt,a4paper]{article}

%used packages
\usepackage{graphicx}
\usepackage[dutch]{babel} 
\usepackage{scrextend}
\usepackage{enumitem}
\usepackage{listings}
\usepackage{float}
\usepackage{titlesec}
\usepackage[table,xcdraw]{xcolor}
\usepackage{fancyhdr}
\usepackage{lastpage}
\usepackage{xcolor}
\usepackage{listings}

%Document settings
\raggedbottom
\graphicspath{ {images/} }
\setlength{\parindent}{0em}
\setlength{\parskip}{1em}

%Custom commando's
\newcommand\litem[1]{\item{\bfseries#1.\space}}
\newcommand\tabelkleur{\rowcolor[HTML]{FFCC67}}

%Custum Variables
\def\auteur{Tim van Broekhoven}
\def\titel{Verslag analoge modulatiemethoden} 
\def\datum{\today}
\def\versie{0.1}
\def\status{concept}
\def\subtitel{}
\def\tabelkleur{FFCC67}

%Settings for Footer and Header
\pagestyle{fancy}
\fancyhf{}
\rhead{\subtitel}
\lhead{\titel}
\lfoot{\datum}
\rfoot{Pagina \thepage \hspace{1pt} /  \pageref{LastPage}}

%Settings for document spacing
\titlespacing{\section}{0pt}{*0}{*0}
\titlespacing{\subsection}{0pt}{*0}{*0}
\titlespacing{\subsubsection}{0pt}{*0}{*0}

%Settings for code sinppeds
\lstset{basicstyle=\ttfamily,
	showstringspaces=false,
	commentstyle=\color{red},
	keywordstyle=\color{blue}
}

\begin{document}
	
	\begin{titlepage}	
		\centering
		{\huge\bfseries \titel \par}
		{\huge\itshape \subtitel \par}
		\vspace{1cm}
		{\Large\itshape \auteur \par}
		\vspace{1cm}
		{\Large\itshape versie \versie\par}	
		\vfill
		\includegraphics[width=0.30\textwidth]{hr-logo}\par\vspace{1cm}
		Vak:\par
		PTE01
		
		\vfill
		{\large \datum \par}
	\end{titlepage}
	
	\tableofcontents
	
	\clearpage
	
	\listoffigures
	
	\clearpage
	\listoftables
	
	\clearpage
	
	\section{Versiehistorie}
	\begin{table}[H]
		\centering
		\label{Versiehistorie}
		\begin{tabular}{|p{1cm}|p{2cm}|p{6cm}|p{2cm}|}
			\hline
			\rowcolor[HTML]{FFCC67}
			\textbf{Versie} & \textbf{Datum} & \textbf{Wijzigingen} & \textbf{Auteur} \\ \hline
			0.1    & 28-01-2018 & Opmaak    & Tim \\ \hline
			0.2	   & ...		& ...  	&  \\ \hline
			&       &             &        \\ \hline
		\end{tabular}
		\caption {Versiehistorie} \label{tab:title} 
	\end{table}	

	\begin{figure}[H]
		\includegraphics[width=1.0\textwidth]{placeimg}\par\vspace{1cm}
		\caption{placeimg}
		\label{fig:length_eight_mouse}
	\end{figure}
	\clearpage
	
	\section{Inleiding}
		\subsection{Probleemstelling}
		\subsection{Doelstelling}
		\subsection{Opdrachtomschrijving}
	\clearpage

%----------------------------------------------------------------------------------------
%	AM Modulatie
%----------------------------------------------------------------------------------------	
	\section{AM Modulatie}
		\subsection{Ontwerpen}
			\subsubsection{Wiskundige afleiding}
			 	\[ S_{AM} = (V_x + m(t)) * c(t) \]

			 \subsubsection{Simulatie}
			 	scoopbeelden + frequentie spectra
		
		\subsection{Meetresultaten} 
			\subsubsection{Meetopstelling}
				Hoe is gemeten?
			\subsubsection{Signaalmetingen}
				(scopebeelden) bij elke signaalbewerking +
frequentiespectrum + tabel vergelijking meetresultaten
		
		\subsection{Conclusie}
			\begin{itemize}
				\item Kloppen de gemeten waarden met de berekende waarden
				\item Waarom wel/niet?
				\item Voor/nadelen van de verschillende schakelingen
				\item Aanbevelingen voor evt. verbeteringen
			\end{itemize}
	\clearpage

%----------------------------------------------------------------------------------------
%	AM Demodulatie
%------------------------------------------------------------------------------------			
	\section{AM Demodulatie}
		\subsection{Ontwerpen}
			\subsubsection{Wiskundige afleiding}
			 	
			 \subsubsection{Simulatie}
			 	scoopbeelden + frequentie spectra
		
		\subsection{Meetresultaten} 
			\subsubsection{Meetopstelling}
				Hoe is gemeten?
			\subsubsection{Signaalmetingen}
				(scopebeelden) bij elke signaalbewerking +
frequentiespectrum + tabel vergelijking meetresultaten
		
		\subsection{Conclusie}
			\begin{itemize}
				\item Kloppen de gemeten waarden met de berekende waarden
				\item Waarom wel/niet?
				\item Voor/nadelen van de verschillende schakelingen
				\item Aanbevelingen voor evt. verbeteringen
			\end{itemize}
	
	\bibliographystyle{plain}
	\bibliography{bibligraphy}

\end{document}